\setlength{\absparsep}{18pt} % ajusta o espaçamento dos parágrafos do resumo

\begin{resumo}[Resumo]
\begin{otherlanguage*}{brazil}

O giro denteado (GD) é responsável pela separação de padrões, função computacional que transforma padrões de entrada similares em
representações neurais mais distintas, e é também uma das únicas áreas do encéfalo de mamíferos que gera novos neurônios ao longo
de toda a vida adulta. As células granulares (CGs) geradas por essa neurogênese passam por uma fase imatura de alta
excitabilidade, e ativá-las \textit{in vivo} recruta interneurônios inibitórios e aumenta a esparsidade da atividade tanto no GD
quanto no Cornu Ammonis 3 (CA3). Ainda não foi investigado, no entanto, se esse efeito sobre o circuito decorre do microcircuito
já conhecido do GD-CA3, visto que os modelos computacionais existentes de neurogênese e separação de padrões se restringem ao
GD\@. Neste trabalho, construiu-se um modelo computacional da rede GD-CA3 a partir da base de conhecimento Hippocampome.org,
estendido com uma população de CGs imaturas (iCGs). A separação de padrões e a esparsidade foram medidas variando-se o nível de
neurogênese e a excitabilidade das iCGs, que em seguida foram ativadas diretamente. A neurogênese aumentou a capacidade de
separação de padrões das CGs maduras (mCGs) e levou o CA3 no mesmo sentido, principalmente por tornar a atividade populacional
mais esparsa; já a ativação direta das iCGs recrutou interneurônios inibitórios por todo o GD e o CA3, aumentando a esparsidade e
reproduzindo o efeito observado \textit{in vivo}. Conclui-se que a influência da neurogênese adulta sobre a esparsidade pode
emergir apenas desse microcircuito e que ela alcança regiões além do próprio GD, corroborando a visão das iCGs como moduladoras da
atividade do circuito, ao invés de codificadoras de informação.

\noindent\textbf{Palavras-chave}: Neurogênese adulta, separação de padrões, giro denteado, CA3, rede neural pulsada, esparsidade.
\end{otherlanguage*}
\end{resumo}

\begin{resumo}[Abstract]
\begin{otherlanguage*}{english}

The dentate gyrus (DG) transforms similar inputs into more distinct representations, a computation termed pattern
separation, and is one of the few brain regions that generate new neurons throughout adulthood. These adult-born granule cells
(GCs) pass through an immature, highly excitable stage, and activating them \textit{in vivo} recruits inhibitory interneurons and
increases population sparsity across both the DG and CA3 (Cornu Ammonis area 3). Whether this circuit-wide effect follows from the
known DG-CA3 microcircuitry has not been examined in a biophysically detailed model, and existing models of neurogenesis and
pattern separation are confined to the DG\@. Here, a spiking model of the DG-CA3 network was built from the Hippocampome.org
knowledge base and extended with a population of immature GCs (iGCs). Pattern separation and population sparsity were measured
while the level of neurogenesis and iGC excitability was varied, and the iGCs were then activated directly. Neurogenesis enhanced
pattern separation in the mature GCs (mGCs) and pushed CA3 in the same direction, mainly by making population activity sparser;
direct activation of the iGCs recruited inhibitory interneurons throughout the DG and CA3 and increased population sparsity,
reproducing the effect observed \textit{in vivo}. The sparsifying influence of adult neurogenesis can therefore emerge from the
known DG-CA3 circuitry and reaches beyond the DG, consistent with a modulatory rather than an encoding role for iGCs.

\noindent \textbf{Keywords}: Adult neurogenesis, pattern separation, dentate gyrus, CA3, spiking neural network, sparsity.
\end{otherlanguage*}
\end{resumo}
